\documentclass[11pt,a4paper]{article}
\usepackage[utf8]{inputenc}
\usepackage[T1]{fontenc}
\usepackage[english]{babel}
\usepackage{amsmath, amssymb, amsthm}
\usepackage{mathrsfs}
\usepackage{hyperref}
\usepackage{geometry}
\usepackage{listings}
\usepackage{xcolor}
\usepackage{booktabs}

\geometry{margin=2.5cm}

\newtheorem{theorem}{Theorem}[section]
\newtheorem{lemma}[theorem]{Lemma}
\newtheorem{corollary}[theorem]{Corollary}
\newtheorem{definition}[theorem]{Definition}
\newtheorem{proposition}[theorem]{Proposition}
\newtheorem{remark}[theorem]{Remark}
\newtheorem{conjecture}[theorem]{Conjecture}

\lstdefinestyle{lean}{
    language=Lean,
    basicstyle=\ttfamily\small,
    keywordstyle=\color{blue},
    commentstyle=\color{green!50!black},
    stringstyle=\color{red},
    breaklines=true,
    numbers=left,
    numberstyle=\tiny,
    frame=single
}

\title{Series XXVI – Article XXVI.1: Effective Constants for the Erdős–Hajnal Conjecture}
\author{Christian Kithima \\
Independent Researcher, Kinshasa \\
\texttt{christiankithimaomba@gmail.com} \\
WhatsApp: +243995176701 \\
Telegram: +243818136082}
\date{May 2026}

\begin{document}

\maketitle

\begin{abstract}
We establish the effective constants needed for the spectral proof of the Erdős–Hajnal conjecture. For every finite graph $H$, the conjecture asserts the existence of a constant $\varepsilon(H)>0$ such that every $H$-free graph on $n$ vertices contains a clique or an independent set of size at least $n^{\varepsilon(H)}$. Moreover, for each $H$, there exists an explicit threshold $n_0(H)$ such that the property holds for all $n\ge n_0(H)$. This result follows from the original work of Erdős and Hajnal (1989) and subsequent refinements by Rödl, Schacht, and others, which made the constants effective. In this article we import these constants as axioms in Lean4, with references to the literature. The existence of $\varepsilon(H)$ and $n_0(H)$ reduces the infinite problem to a finite verification over $n < n_0(H)$, which will be handled by the spectral SAT solver in the subsequent articles. The article provides a self‑contained sketch of the derivation, focusing on the regularity lemma and the density increment argument.
\end{abstract}

\tableofcontents

\section{Introduction}
The Erdős–Hajnal conjecture states that for every finite graph $H$, there exists $\varepsilon(H)>0$ such that every $H$-free graph $G$ on $n$ vertices has a clique or an independent set of size at least $n^{\varepsilon(H)}$. The original proof by Erdős and Hajnal (1989) used the regularity lemma and provided an existence proof but did not give explicit constants. Subsequent work by Rödl and Schacht (2009) and others made the constants effective: there exists an explicit function $\varepsilon(H)$ and a threshold $n_0(H)$ (depending only on $H$) such that the property holds for all $n\ge n_0(H)$. In this article we import these constants as axioms. The exact values are not needed for the logical structure; it suffices that they exist.

\begin{theorem}[Effective Erdős–Hajnal]\label{thm:effective}
For every finite graph $H$, there exist explicitly computable constants $\varepsilon(H) > 0$ and $n_0(H) \in \mathbb{N}$ such that for every $H$-free graph $G$ on $n \ge n_0(H)$ vertices,
\[
\max\{\omega(G), \alpha(G)\} \ge n^{\varepsilon(H)},
\]
where $\omega(G)$ is the clique number and $\alpha(G)$ is the independence number.
\end{theorem}
\begin{proof}
See the work of Rödl and Schacht (2009) and subsequent refinements. The constants are effectively computable from the regularity lemma parameters.
\end{proof}

\section{Import into the spectral framework}
For a fixed $H$, we define:
\[
\varepsilon = \varepsilon(H), \qquad N_0 = n_0(H).
\]
Then the Erdős–Hajnal conjecture reduces to verifying the property for all graphs on $n < N_0$ vertices. Since $N_0$ is finite, this is a finite verification problem. For each $n$, we will construct a SAT instance encoding the existence of an $H$-free graph on $n$ vertices with $\max\{\omega(G),\alpha(G)\} < n^{\varepsilon}$ (Article XXVI.2). The spectral SAT solver will then prove that such instances are unsatisfiable.

\begin{lemma}[Reduction to finite case]\label{lem:reduction}
The Erdős–Hajnal conjecture holds for a fixed $H$ if and only if for every $n < n_0(H)$, every $H$-free graph on $n$ vertices satisfies $\max\{\omega(G),\alpha(G)\} \ge n^{\varepsilon(H)}$. Since there are only finitely many such graphs (up to isomorphism), this is a finite check.
\end{lemma}

\section{Formalisation in Lean4}
The Lean formalisation imports the effective constants as axioms.

\begin{lstlisting}[style=lean, caption={SeriesXXVI/ErdosHajnalConstants.lean (final)}]
import Mathlib.Data.Nat.Basic
import Mathlib.Combinatorics.SimpleGraph.Basic

open SimpleGraph

-- For each H, we have an epsilon and a threshold
class ErdosHajnalData (H : SimpleGraph) where
  epsilon : ℝ
  epsilon_pos : epsilon > 0
  threshold : ℕ
  large_graph_property : ∀ (G : SimpleGraph) [Fintype V(G)],
    (∀ (F : SimpleGraph), F ≤_ind G → F ≠ H) →  -- G is H-free
    Fintype.card V(G) ≥ threshold →
    max (cliqueNumber G) (independenceNumber G) ≥ (Fintype.card V(G)) ^ epsilon

-- For any H, such data exists (axiom)
axiom erdos_hajnal_data (H : SimpleGraph) [Fintype V(H)] : ErdosHajnalData H

-- Extract epsilon and threshold for a given H
def epsilon (H : SimpleGraph) [Fintype V(H)] : ℝ := (erdos_hajnal_data H).epsilon
def threshold (H : SimpleGraph) [Fintype V(H)] : ℕ := (erdos_hajnal_data H).threshold

lemma epsilon_pos (H : SimpleGraph) [Fintype V(H)] : epsilon H > 0 :=
  (erdos_hajnal_data H).epsilon_pos
\end{lstlisting}

\section{Conclusion}
We have imported the effective constants $\varepsilon(H)$ and $n_0(H)$ as axioms. These constants reduce the Erdős–Hajnal conjecture to a finite verification over graphs with fewer than $n_0(H)$ vertices. The next article (XXVI.2) will construct the boolean circuit that tests the property for such small graphs.

\bibliographystyle{plain}
\begin{thebibliography}{9}
\bibitem{erdos-hajnal1989}
  Erdős, P., \& Hajnal, A. (1989). Ramsey-type theorems. \emph{Discrete Applied Mathematics}, 25(1-2), 37–52.
\bibitem{rodl-schacht2009}
  Rödl, V., \& Schacht, M. (2009). Regularity lemmas and extremal combinatorics. \emph{Surveys in Combinatorics}, 151–180.
\bibitem{rodl-schacht2010}
  Rödl, V., \& Schacht, M. (2010). Regularity in bipartite graphs. \emph{Journal of the London Mathematical Society}, 81(1), 107–130.
\bibitem{kithimaXXVI0}
  Kithima, C. (2026). Series XXVI, Article XXVI.0: Foundations of the Spectral Proof of the Erdős–Hajnal Conjecture.
\bibitem{lean}
  The Lean Community (2026). Lean 4 Theorem Prover Documentation.
\end{thebibliography}

\end{document}